\chapter{Introduction}
\label{chap:introduction}

\begin{drafttext}
Economics and AI alignment begin from the same structural problem: how can a
system produce coherent, desirable behaviour when relevant information,
capability, and control are distributed among actors with different objectives?
Economic theory studies institutions that turn local preferences and private
information into prices, allocations, and incentives. AI alignment studies how
to make increasingly capable artificial agents pursue objectives that remain
desirable under novelty, scale, and informational advantage. The analogy is
therefore not merely metaphorical. Markets provide both examples of composite
agency and a mathematical language for reward, information, and delegation.
\end{drafttext}

Before neural networks won the mandate of heaven, an AI research paradigm that
had shown considerable promise was that of \emph{market-based architectures},
i.e. AI agents that determine their output based on some internal market-based
mechanism \citep{kweeMarketBasedReinforcementLearning2001,
changDecentralizedReinforcementLearning2020}.

More broadly speaking, there are general intuitive arguments to imagine a close
analogy between markets and AI. Philosophers and psychologists have long
pondered multi-agent models of the mind \citep{minskySocietyMind1988}; moreover,
one may imagine that ``any'' machine learning task could in principle be solved
by a market of agents solving sub-tasks with their individual reward set by the
sale value of their output. There is work suggesting that markets can capture
some notion of \emph{bounded rationality}, e.g. the \emph{Boundedly Rational
Inductive Agent} (BRIA) \citep{oesterheldTheoryBoundedInductive2023} and
\emph{Algorithmic Bayesian Epistemology}
\citep{neymanAlgorithmicBayesianEpistemology2024}. In some sense, markets
``aggregate'' the intelligence or capacities of their individual participants.

\section{Three levels of the economic alignment problem}

\begin{drafttext}
The thesis is organized around three levels. At the first level, economic
mechanisms are studied \emph{as agents}: can market interaction among bounded
sub-agents implement reinforcement learning, backpropagation, or rational reward?
At the second level, economics is used to study the relationship between an
agent and its evaluator: scalable oversight becomes an information-asymmetry
problem in which the party being evaluated may understand its output better
than the party assigning reward. At the third level, the object to be elicited
is belief itself: how can a mechanism score forecasts or propositions when
ground truth is delayed, non-verifiable, non-falsifiable, or partly subjective?

The third level connects to the first two through information economics. Prices
can guide agency only to the extent that they encode useful beliefs, while an
oversight mechanism can reward an agent only to the extent that the evaluator
can elicit or test those beliefs. Prediction markets, proper scoring rules,
arbitrage, and consistency checks are therefore epistemic infrastructure for the
same economic account of aligned agency.
\end{drafttext}

\section{Contributions and structure}

\begin{drafttext}
Part~I studies market-based reinforcement learning and a price-recursion account
of reward. Part~II develops the information-asymmetry framing of scalable
oversight, recursive information markets, and a benchmark for comparing
oversight mechanisms. Part~III studies belief elicitation beyond promptly
resolvable questions, first through verification--falsification games and then
through logical consistency and arbitrage among forecasts.
\end{drafttext}
